\section*{Capítulo \ref{ch:b2:muscle-movement} --- Contracción muscular y función motora}

\begin{solution}{exo:b2:muscle-movement:1}
Los filamentos finos se deslizan a lo largo de los gruesos, arrastrados por
las cabezas de miosina; ningún filamento se acorta. La banda I y la zona H
se estrechan y los discos Z se acercan; la banda A (los filamentos gruesos)
conserva su longitud.
\end{solution}

\begin{solution}{exo:b2:muscle-movement:2}
(1) ATP unido, cabeza separada; ATP hidrolizado, cabeza armada. (2) La
cabeza se une a la actina. (3) Liberación del fosfato, golpe de fuerza,
liberación del ADP. (4) Un nuevo ATP se une y la cabeza se separa. El ATP
es necesario para \emph{soltar} la cabeza: sin él, todas las cabezas
quedan unidas y el músculo se bloquea --- rigidez.
\end{solution}

\begin{solution}{exo:b2:muscle-movement:3}
Impulso nervioso; liberación de acetilcolina y potencial de placa terminal
(\qty{0.6}{ms}); \omterm{prop:b2:neurons-synapses:ap}{potencial de acción} muscular a lo largo de la fibra y por
los \omterm{def:b2:cell-differentiation:fibre}{túbulos T} (\qty{1}{ms}); liberación de calcio desde el \omterm{def:b2:cell-differentiation:fibre}{retículo
sarcoplásmico} (\qty{1}{ms}); el calcio se une a la troponina, la
tropomiosina se desplaza, las cabezas se unen (unos pocos milisegundos);
la fuerza sube a lo largo de decenas de milisegundos.
\end{solution}

\begin{solution}{exo:b2:muscle-movement:4}
\omterm{def:b2:muscle-movement:motorunit}{Unidad motora}: una motoneurona y todas las fibras que inerva. Sacudida: la
contracción breve producida por un solo impulso. \omterm{def:b2:muscle-movement:motorunit}{Tétanos}: la contracción
fusionada y sostenida producida por impulsos que llegan más deprisa de lo
que decae la sacudida. La fuerza se gradúa mediante la frecuencia de
disparo de cada unidad y mediante el número de unidades reclutadas.
\end{solution}

\begin{solution}{exo:b2:muscle-movement:5}
Cada hemisarcómero se desliza \qty{0.2}{\micro m} en \qty{50}{ms}:
\qty{4}{\micro m/s}; son 20 golpes de \qty{10}{nm} en \qty{50}{ms}, 400 por
segundo si una cabeza estuviera unida todo el tiempo.
\end{solution}

\begin{solution}{exo:b2:muscle-movement:6}
\qty{2}{cm} en \qty{0.1}{s}: \qty{0.2}{m/s}, 2 longitudes por segundo;
cada \omterm{def:b2:cell-differentiation:fibre}{sarcómero} se acorta $2/\num{40000}$ cm $= \qty{0.5}{\micro m}$ en
\qty{0.1}{s}, \qty{5}{\micro m/s}.
\end{solution}

\begin{solution}{exo:b2:muscle-movement:7}
$F/F_0 = 0.3125/(v/v_{\max} + 0.25) - 0.25$: a 1, 3 y 6 longitudes por
segundo ($0.1$, $0.3$, $0.6$ de $v_{\max}$): $0.64$, $0.32$, $0.12$.
Potencia $Fv$ (en $F_0\times$ longitudes por segundo): $0.64$, $0.95$, $0.71$
--- máxima hacia 3 longitudes por segundo, alrededor de un tercio de
$v_{\max}$.
\end{solution}

\begin{solution}{exo:b2:muscle-movement:8}
$80\times 30 = \qty{2400}{N}$; en el pie, $2400\times 5/45 =
\qty{270}{N}$.
\end{solution}

\begin{solution}{exo:b2:muscle-movement:9}
A \qty{10}{Hz}, los impulsos llegan cada \qty{100}{ms}, cuando el calcio
del anterior ya ha desaparecido: sacudidas separadas, cada una partiendo
de un estado parcialmente relajado, que dan una ondulación. A
\qty{50}{Hz} llegan cada \qty{20}{ms}, antes de que el calcio se haya
bombeado: el calcio sigue alto y la fuerza se fusiona. El \omterm{def:b2:muscle-movement:motorunit}{tétanos} supera
la sacudida porque un solo pulso de calcio es demasiado breve para que los
puentes cruzados estiren los elementos elásticos en serie y desarrollen
toda la fuerza; el calcio sostenido se lo permite.
\end{solution}

\begin{solution}{exo:b2:muscle-movement:10}
Trabajo, \qty{20}{kJ}; energía del ATP, \qty{80}{kJ}; $1.6$ mol de ATP.
Primeros \qty{5}{s}: $0.8$ mol de fosfocreatina. Los $0.8$ mol de ATP
restantes, por glucólisis: $0.4$ mol de glucosa, $0.8$ mol de lactato ---
en \qty{40}{L}, \qty{20}{mmol/L}.
\end{solution}

\begin{solution}{exo:b2:muscle-movement:11}
Canal de liberación bloqueado: sin calcio no hay contracción --- parálisis.
Bomba bloqueada: el calcio sigue alto y el músculo no puede relajarse ---
una contractura. Troponina insensible: la tropomiosina nunca se desplaza,
no hay contracción. Sin calcio en la \omterm{def:b2:blood-circulation:blood}{sangre}: el músculo esquelético se
contrae con normalidad (su calcio es interno); el \omterm{def:b2:heart:anatomy}{corazón}, que necesita
calcio externo para desencadenar su liberación, se detiene.
\end{solution}

\begin{solution}{exo:b2:muscle-movement:12}
La cabeza convierte en trabajo hasta la mitad de la energía libre de un
ATP; el resto es calor, y las pérdidas de las bombas, de los puentes
cruzados que se unen sin tirar a gran velocidad y de los elementos
elásticos reducen el conjunto a una cuarta parte --- el rendimiento de un
buen motor de gasolina. Las dos maneras de graduar la fuerza permiten a la
vez un control fino en los esfuerzos pequeños (unidades pequeñas añadidas
de una en una, cada una disparando más deprisa o más despacio) y un gran
intervalo (de un factor mil, desde una sola unidad pequeña a frecuencia de
sacudida hasta todas las unidades en \omterm{def:b2:muscle-movement:motorunit}{tétanos}); un simple interruptor daría
una fuerza o ninguna.
\end{solution}

\begin{solution}{pb:b2:muscle-movement:1}
\textbf{1.} $v = \sqrt{2\times 9.81\times 0.4} = \qty{2.8}{m/s}$;
$\tfrac12\times 70\times 2.8^{2} = \qty{275}{J}$.
\textbf{2.} Trabajo $= 275 + 70\times 9.81\times 0.4 = \qty{550}{J}$ sobre
\qty{0.4}{m}: \qty{1370}{N}, el doble del peso.
\textbf{3.} $F_0 = 500\times 30 = \qty{15000}{N}$; en el pie,
\qty{1500}{N}.
\textbf{4.} \qty{4}{cm} en \qty{0.3}{s}: \qty{0.13}{m/s}, $1.1$ longitudes
por segundo, $0.14\,v_{\max}$.
\textbf{5.} $F/F_0 = 0.3125/(0.14 + 0.25) - 0.25 = 0.55$: \qty{8300}{N}
en el músculo, \qty{830}{N} en el pie --- muy lejos de los
\qty{1370}{N} necesarios. El acortamiento muscular por sí solo no puede
realizar este salto. El contramovimiento aporta el resto: cuando el
saltador se agacha, los tendones estirados almacenan energía elástica, y el
músculo, contrayéndose de forma casi isométrica con una fuerza elevada, los
carga; después los tendones se retraen más deprisa de lo que puede
acortarse ninguna fibra.
\textbf{6.} $550/0.3 = \qty{1.8}{kW}$; \qty{92}{W/kg} de extensor.
\textbf{7.} $4/\num{48000}$ cm $= \qty{0.83}{\micro m}$ por \omterm{def:b2:cell-differentiation:fibre}{sarcómero},
\qty{0.42}{\micro m} por hemisarcómero: 42 golpes de \qty{10}{nm}.
\textbf{8.} \qty{0.42}{\micro m} en \qty{0.3}{s}: \qty{1.4}{\micro m/s};
140 golpes por segundo si estuviera unida de forma continua.
\textbf{9.} $500\times 6\times 10^{10}\times \num{48000}\times 300 =
4.3\times 10^{20}$ cabezas.
\textbf{10.} $0.55\times \num{15000}\times 0.04 = \qty{330}{J}$; por
golpe, $0.4\times 8\times 10^{-20} = \qty{3.2e-20}{J}$; $1.0\times
10^{22}$ golpes.
\textbf{11.} $1.0\times 10^{22}/4.3\times 10^{20} = 24$ golpes por
cabeza, de los 42 disponibles: alrededor del \qty{60}{\%}. La reserva es
pequeña --- el músculo trabaja cerca de su límite, como mostró la
pregunta 5.
\textbf{12.} $1.0\times 10^{22}/6.0\times 10^{23} = 0.017$ mol,
\qty{8.7}{g}; su energía, $0.017\times 50 = \qty{860}{J}$, de los que
\qty{330}{J} se convirtieron en trabajo: \qty{39}{\%}.
\textbf{13.} Más de un impulso cada \qty{30}{ms}: por encima de
\qty{33}{Hz}; unos 10 impulsos en \qty{0.3}{s}.
\textbf{14.} $9.9\times 10^{-6}\times 5\times 10^{-10}\times 6\times
10^{23} = 3\times 10^{9}$ iones libres; $1.5\times 10^{9}$ ATP.
\textbf{15.} Fibras: $500/(\pi\times(3\times 10^{-3})^{2}) = 1.8\times
10^{7}$; ATP de los puentes cruzados por fibra y por impulso: $1.0\times
10^{22}/(1.8\times 10^{7}\times 10) = 6\times 10^{13}$, cuarenta mil veces
la estimación de la bomba. El calcio libre es un pequeño residuo: el
retículo libera del orden de \qty{0.1}{mmol/L}, diez veces la subida del
calcio libre, que la troponina y los tampones fijan casi por completo al
instante --- y todo él debe bombearse de vuelta.
\textbf{16.} Para un salto máximo, casi todas las unidades en unas pocas
decenas de milisegundos, primero las pequeñas y al final las grandes; un
paso suave recluta quizá entre una décima y una quinta parte de ellas.
\textbf{17.} Las unidades se reclutan por orden de tamaño: en los
esfuerzos pequeños, cada nueva unidad añade un pequeño incremento; en los
grandes, las unidades restantes son grandes y cada una añade un escalón
importante.
\textbf{18.} El potencial de placa terminal de todas las fibras se reduce
a la mitad; las fibras cuyo potencial cae por debajo del umbral no
disparan, de modo que la sacudida y el \omterm{def:b2:muscle-movement:motorunit}{tétanos} son menores, y el salto se
queda corto.
\textbf{19.} $5\times 20 = \qty{0.1}{mol}$ de ATP: lo que cuestan seis
saltos.
\textbf{20.} $20\times 20 = \qty{0.4}{mol}$ de fosfocreatina: unos 23
saltos.
\textbf{21.} Diez saltos consumen $0.17$ mol, sin agotar todavía la
fosfocreatina; tras una veintena, la glucólisis toma el relevo, y se
acumulan lactato y protones.
\textbf{22.} \qty{1.8}{kW} de trabajo al \qty{40}{\%} son \qty{4.6}{kW} de
ATP, $0.09$ mol por segundo, que necesitan $0.015$ mol de oxígeno por
segundo --- \qty{0.34}{L/s}, \qty{21}{L/min}, diez veces lo que puede
aportar la \omterm{def:b2:blood-circulation:blood}{sangre}: el salto es anaerobio por necesidad, y el oxígeno llega
después.
\textbf{23.} $0.17$ mol de ATP necesitan $0.029$ mol de oxígeno,
\qty{0.64}{L}; pagados en un minuto, \qty{640}{mL/min}, dos veces y media
el consumo en reposo.
\textbf{24.} El ATP se refosforila a partir de la fosfocreatina en
milisegundos, de modo que su concentración apenas baja; la rigidez exige
un ATP cercano a cero, que un músculo vivo con alguna reserva nunca
alcanza.
\textbf{25.} 42 golpes disponibles por hemisarcómero, 24 necesarios;
$F/F_0 = 0.55$ a la velocidad del salto; $0.017$ mol, \qty{8.7}{g}, de ATP.
\end{solution}
